\section{Benchmark Edge Runtime}

Benchmark Edge Runtime đo toàn pipeline trong container AI Runtime thay vì chỉ
engine TensorRT. Các chỉ số được lấy trực tiếp từ telemetry \texttt{runtime} mà
Runtime gửi về nền tảng (Mục~\ref{sec:monitoring-design}), nên phép đo này đồng
thời kiểm chứng luồng telemetry. Nguồn vào là tệp video công trường phát ở 15 FPS
(\texttt{target\_fps}=15).

\begin{table}[H]
\centering
\caption{Chỉ số pipeline AI Runtime lấy từ telemetry (cấu hình Pruned-30 + distill)}
\label{tab:runtime-pipeline-benchmark}
\begin{tabularx}{\textwidth}{L{6cm}YY}
\toprule
\textbf{Chỉ số telemetry} & \textbf{Jetson Nano (FP16)} & \textbf{Orin Nano (INT8)}\\
\midrule
\texttt{inference\_fps} & \cbs{} & \cbs{}\\
\texttt{effective\_fps\_ratio} & \cbs{} & \cbs{}\\
\texttt{inference\_p50\_ms} / \texttt{p95\_ms} & \cbs{} & \cbs{}\\
\texttt{camera\_frame\_drop\_pct} & \cbs{} & \cbs{}\\
\texttt{runtime\_queue\_size} (trung bình) & \cbs{} & \cbs{}\\
\bottomrule
\end{tabularx}
\end{table}

Vì hàng đợi của Runtime giới hạn 8 khung và loại khung cũ khi đầy, latency không
tăng tích lũy khi detector chậm hơn camera; thay vào đó tỷ lệ frame drop tăng.
Do đó \texttt{effective\_fps\_ratio} và \texttt{camera\_frame\_drop\_pct} là hai chỉ
số trực tiếp cho biết cấu hình có đáp ứng thời gian thực hay không.
